
\question[4]
Erkläre, was ein Debugger ist und wozu er beim Programmieren genutzt wird. 
Gehe dabei auf die Begriffe \emph{Einzelschritt}, \emph{Breakpoint} und \emph{Fortfahren} ein.

\begin{solutionbox}{6cm}
Ein Debugger ist ein Werkzeug der Entwicklungsumgebung, mit dem
man ein Programm schrittweise (Anweisung für Anweisung) ausführen
und dabei die aktuellen Werte der Variablen beobachten kann.
So lassen sich logische Fehler leichter finden, da man genau
sieht, wo sich ein Wert anders verändert als erwartet.

Einzelschritt: Führt jeweils nur die nächste Anweisung aus.

Breakpoint: Eine markierte Zeile, an der das Programm automatisch
anhält, sodass man nicht jede Zeile einzeln durchgehen muss.

Fortfahren: Lässt das Programm laufen, bis der nächste Breakpoint
erreicht wird (oder das Programm endet).
\end{solutionbox}